\section{Applying the local estimates to one marked step}
\label{sec:actual-local-compiler-handoff}

The local estimates must hold for the very matrices and profiles used in
the marked density identity. We now make this common choice explicit.
The argument combines the preceding local results; it does not introduce
a new Gaussian law or replace the selected history.

Fix a gauge group, a patch of block size $L$, its funded tube construction,
and the probability-Haar chart of Section~\ref{sec:haar-chart}. The
minimizing and fallback branches are those fixed in
Section~\ref{sec:predictor-branch}. Let $e\geq0$ be the support exponent
needed by the marked step, including its prescribed reserve, and let
$t>0$ be its footprint allowance. The conversion from fine supports to
coarse connected covers is required to have reserve $t$ at this $e$.
Fix also the profile width $W>0$, the chart-radius multiplier $d\geq0$,
the warning margin $\varepsilon>0$, the owner-energy constant $m>0$,
and $K>0$ with
\[
 2K\leq m/2,\qquad 2a_{\rm coast}\leq K/4.
\]
Here $a_{\rm coast}$ is the incoming coefficient in the small-remainder
budget. These are the numerical margins already used by the corrected
compiler and its inherited-activity estimate.

The guide estimate is chosen before $L$: the absolute path sums of
Section~\ref{sec:seed-guides} give $C=2A(1+a)$ simultaneously for all
coarse fields, including all stores later compatible with those fields.
The choices in Section~\ref{sec:warning-safety} then impose
\[
 C/L^2\leq\sqrt{\mathrm{direct}}/64,
 \qquad C+\sqrt{\mathrm{direct}}/64
              \leq (W/4)\sqrt{\mathrm{penalty}}.
\]
After $L$ and $W$ have been fixed, the remaining constants below may depend
on this patch. They are chosen before the coupling parameter and before
any torus, forest, tag, coarse field or finite family of insertions.

For a large parameter $x$, write $y=r(x)$ for its retiming and
$B=\lceil(\log y)^2\rceil$. An admissible marked step consists of its
actual tube, fine representation $R$, coarse budget, chart and profiles.
Its tube is funded by the fixed patch construction; its upward paths and
high clips are the path and clip of $R$. Its coarse threshold is
$\operatorname{temp}(y)$, its chart radius is $d\,\mathrm{fiery}(y)$,
and its coarse side is at least $100(1+B)$. The remaining radius,
threshold, coupling-band and inherited-law inequalities are exactly those
listed in Section~\ref{sec:actual-corrected-body-compiler}. In addition we
retain the actual alongside certificate for the tube and the small
additive remainder of $R$. These identify the input to the construction;
large coupling alone would not identify its paths or its remainder.

\subsection{One matrix certificate and both support inventories}

For every compatible store, forest, original tag and coarse field, use the
actual full precision $A$ of Section~\ref{sec:sector-gram}. Its selected
Gram matrix and its selected screened covariance have positive floors.
The bounds are uniform over all of these data, including frozen
coordinates; compatibility requires that the stored coarse flags equal
the flags of the coarse field at threshold $\operatorname{temp}(y)$.

Here is a simultaneous choice of constants in the local inverse and
logarithmic estimates. Let $s>0,M_0\geq0,\eta>0$ be their common initial
Schur rate, amplitude and gap. Set $M=\max(M_0,1)$ and
$h=\min(\eta,1)$. Choose the inverse rate $Y>0$ supplied by
Section~\ref{sec:local-covariance}, with $Y<s$, and put
\[
 m_0=Y/2,\quad v=m_0/2,\quad
 R_0=(m_0-v)/120,\quad b=R_0/2,\quad \kappa=b/2.
\]
Enlarge the common amplitude $C\geq1$ to cover the precision, its inverse
and the local volume sum at rate $Y-m_0$, divided by
$\min(h/2,1)$. These constants depend only on the fixed patch data.
The window formulas of Section~\ref{sec:local-covariance} now price the
actual inverse and both determinant corrections with rate $v$, amplitude
$C$ and error $e^{-bB}$. They apply also to every admissible recut whose
windows contain their radius-$B/60$ neighborhoods. Indeed, the missed
entries have separation at least $\lfloor B/60\rfloor$; for $B\geq120$,
\[
 B\leq120\lfloor B/60\rfloor.
\]
The weighted inverse estimate therefore bounds their error by
$C^4e^{-R_0B}\leq e^{-bB}$ eventually. The same window estimate gives
positivity of the recut covariance. The original selected covariance is
kept fixed when these auxiliary recuts are considered.

We also need rooted prices for two actual word inventories, of orders
$n=0$ and $n=1$, with halo multiplier $l=20$. Set
$\theta=b/(4(l+2))$. In the itinerary summation of
Section~\ref{sec:itinerary-prices}, the summation ratio has the form
$C_1e^{-(b-\theta)B}$ and the numerator for order $n$ is bounded by
\[
 C_2e^{-[b-\theta(n+1+l)]B}.
\]
The constants include the fixed cover weight at exponent $e$. Since
$n\leq1$, the exponent in brackets is at least $3b/4>\kappa$.
The ratio is eventually at most $1/2$, and the geometric series is
therefore at most $2C_2e^{-3bB/4}\leq e^{-\kappa B}$.
Intersect the two resulting tails. We obtain both rooted inventory
prices with the same $v,C,b,\kappa$, the actual physical address map,
and the complete halo supports. This is a price for the individual
ordered words through a fixed site, not just for their sum as an operator.

\subsection{The first derivative of the two long determinant families}

The value-level compiler also needs first derivatives as the coarse field
moves in its local logarithmic coordinates. Freeze stores, forests, tags
and windows and use the actual matrix $A(U)$, rather than selecting a new
sector at each nearby field. We first justify a radius independent of the torus. The finite-patch
predictor estimate of Section~\ref{sec:predictor-branch} gives a common
chart radius for its retained rows. For the path rows, use the path
extension bound in Section~\ref{sec:ex-row-quadratic-price} before summing
the slots. A fixed number of derivatives costs a polynomial in the path
length; every additional spatial or derivative-coordinate weight costs
an exponential in that length with an arbitrarily small fixed rate.
Choose the latter rate so that $(1+2j)\delta$ is at most half the
incoming path rate for the required order $j=2$. The incoming absolute
path budget absorbs the polynomial and all these weights. Summing the
absolute slot bounds, including repeated links, therefore gives one
entrywise derivative majorant for all rows on a fixed logarithmic ball.
The local predictor rows have the same property on a smaller ball.
We retain the derivative-coordinate weights in this statement. For an
ordered list $\alpha$ of atomic link-color derivatives, $|\alpha|\leq2$,
put $W_\alpha(o)=\prod_{a\in\alpha}e^{\delta d(o,\mathrm{hearth}(a))}$.
For a row anchored at $o$ and a column anchored at $z$, the bounds are
\[
 W_\alpha(o)e^{\delta d(o,z)}
            |\partial_\alpha M_{oz}|\leq C_1,\qquad
 W_\alpha(z)e^{\delta d(o,z)}
            |\partial_\alpha M_{oz}|\leq C_1.
\]
To verify them term by term, a derivative is zero if any atom misses the
slot's complete support. Otherwise its endpoint and derivative-coordinate
distances cost at most $(1+2|\alpha|)\delta(\ell+5)$ for a slot of
length $\ell$. The chosen half-rate reserve pays this exponential, while
the other half pays its fixed-order differentiation factor. Absolute
summation over the original path labels proves both displayed bounds.
The predictor contribution has fixed reach and a uniform analytic
radius, so the same argument uses its finite support instead. Repeated
atomic derivatives and repeated physical links are included.

Sum rows and columns at spatial rate $\delta/2$, leaving a decaying
factor for their exponential shells. The bounded number of coordinates
per site makes these sums finite uniformly in the period. For derivative
order $k=2$, lower both spatial and derivative-coordinate rates further
to $u=(\delta/2)/(1+k)$. The inequality $u+ku\leq\delta/2$ pays the
transfer of derivative-coordinate weights across the intermediate matrix
index. For $M^{\mathsf T}M$, expand each
ordered derivative list over its finitely many Leibniz partitions. The
triangle inequality transfers the derivative-coordinate weights to the
appropriate factor in each term; weighted row/column multiplication then
bounds that term. Summing the finitely many partitions gives the full
weighted derivative prices for the precision. Complete inactive
coordinates by the identity. All selections remain frozen; no derivative
of a flag indicator occurs. Thus the estimates retain the atomic-list
and support information required by the itinerary argument below.

For precision, these first-derivative bounds give constants
$\sigma,C_0,r_1>0$, uniform in the period and base field, such that
\[
 |(DA(H)Z)_{ij}|\leq C_0S_G(\sigma/2)
                  e^{-\sigma d(i,j)}\|Z\|_\infty
       \qquad(\|H\|_\infty\leq r_1).
\]
Here $S_G(\sigma/2)$ is a uniform upper bound for the sum of the
decaying reciprocal derivative-coordinate weights
$e^{-\sigma d(o,\mathrm{hearth}(a))/2}$, including color multiplicity. Let $V_0$ bound the row and column
volume sums at rate $\sigma/2$, and set
$U_0=V_0C_0S_G(\sigma/2)$. Integrate the last estimate along the segment
from $0$ to $H$ and sum the remaining exponential weights. The Schur
norm of $A(H)-A(0)$ is at most $U_0\|H\|_\infty$.
If the base gap is $\eta$, choose
\[
 r_0=\min\{r_1,\eta/[2(1+U_0)]\}.
\]
Then for every real vector $z$ and $\|H\|_\infty\leq r_0$,
\[
 z^{\mathsf T}A(H)z
 \geq(\eta-U_0\|H\|_\infty)|z|^2
 \geq(\eta/2)|z|^2.
\]
This is a uniform neighborhood gap, including fields that no longer obey
the original flag condition. The flag condition was used at the base
field, and the matrix formula is extended with its selections fixed.

Apply the inverse and window derivative formulas of
Section~\ref{sec:normalization-jets} on this ball. Restrict once more to
half its radius. The actual selected covariance agrees there, as a germ,
with the symmetric window formula for the completed precision: the
unused-coordinate blocks are identities and the active blocks are the
original selected ones. Both this covariance and the precision have the
positive floors supplied by the window estimates. Restrict the logarithmic
direction further to
\[
 \|H\|_\infty\leq\min\{r_0/2,r_G(1)/2\},
\]
where $r_G(1)>0$ is the gauge-coordinate radius at order one. This radius
concerns the variation of a particular matrix function. It is not the
width of the replacement profiles.

Apply the differentiated word estimates of
Section~\ref{sec:itinerary-prices} at derivative order one, halo $20$ and
support exponent $e$. They cover both long determinant classes on the
same field neighborhood. Choose their positive decay rate $d_0$ and
fixed letter factor $F_0$ before $y$. The scalar conversion and the two
classes together cost at most
\[
 4F_0 e^1\{2(1+1)\}^{1}e^{-d_0B}.
\]
For large $B$ this is at most $e^{-d_0B/2}$. Thus the common neighborhood,
individual word prices, and total connected-cover price are available
simultaneously. There is no need to identify $d_0$ with the rate
$\kappa$ of the preceding value inventories; their two positive rates
can be retained separately. The order parameters used here are: scalar
conversion order two, coarse-field derivative order one, and halo twenty.

\subsection{Uniformity on a whole coupling interval}

The palette estimate of Section~\ref{sec:actual-palette-warning} provides
one bound $\ell\geq1$ for all actual pilots of this patch. It gives the
Gaussian-coordinate bound $y^2$ on the door event, the linear-residual
bound $y^{-4}$, and the all-compatible-palette zero-socket conclusion.
The actual owner-energy estimate
\eqref{eq:owner-actual-heft} supplies the other missing input to
Section~\ref{sec:actual-corrected-body-compiler}. Its path counting and
clip estimate use precisely the path and clip identified above. The
fixed chart and branch construction, the guide reserve, and the small
span estimate supply the remaining local premises of these applications.

These assertions must hold for every $x\in[s/2,s]$ once the chosen
station parameter $s$ is large. There are only finitely many eventual
predicates at this stage: the matrix certificate and its two inventories,
the first-derivative word prices, palette bound, owner estimate, local
chart and room conditions, and the fixed numerical inequalities for the
corrected alphabet. Pull the predicates in $y$ back along $r(x)\to\infty$.
Intersect them with the predicates already expressed in $x$, and choose
one lower threshold $x_0$. Finally require $s/2\geq x_0$.
Then every $x\in[s/2,s]$ satisfies them. This step is performed before
quantifying over periods, fields, forests and tags because each local
estimate above was uniform in those variables. It is also before the
choice of any finite marked family.

\subsection{The genuine density and its auxiliary prices}

Apply the corrected exponent decomposition and bounded extension of
Section~\ref{sec:actual-corrected-body-compiler} with these actual data.
Apply the mixed compiler of
Section~\ref{sec:actual-mixed-marked-compiler} using the same word
inventories and full read supports. The resulting profile system and
finite unmarked alphabet are chosen before an output representation
$Q$ with the prescribed coarse budget, and before the marked insertions.

There are two distinct uses of the source representation. Removing its
additive remainder gives $R^\circ$, which supplies the inherited-law,
warning and price estimates. The small-remainder estimate of $R$ itself
is supplied separately. The correction construction adds this genuine
remainder to the alphabet and pays it using $2a_{\rm coast}\leq K/4$.
Thus the planted identity of Section~\ref{sec:allarity-planted-compiler}
reconstructs the density of $R$, whereas the connected-support and return
estimates use $R^\circ$. Both assertions concern exactly the same
corrected alphabet, profiles, path, clip and scalar loading.

That loading is positive: its chart scale and Gaussian prefactor are
positive, the fine coupling is positive, and its reference kinetic Gram
matrix has positive determinant by the common matrix certificate.
The square-root determinant in its denominator is therefore nonzero.
It is the original reference-forest/tag loading, not an independently
normalized version of the corrected alphabet.

At a short interval of the selected history, identify its actual patch
and funded tube using the constant patch data on that interval. The
inherited field at its upper end supplies $R$ and its remainder estimate;
its neutral version supplies the admissible local data above. The
interval length fixes the permitted remainder budget, and the station
bounds put its coupling in $[s/2,s]$. Hence the same construction applies
to each step of that interval.

The restart construction of Section~\ref{sec:restarted-rank-history}
sets the new coarse coupling equal to the incoming one. Use the positive
loading to choose its scalar vacuum normalization, and use the connected
expansion of the same alphabet to define its masked logarithmic
remainder. This yields the exact pushed density and preserves the
original coarse path and clip. It does not assert equality with the old
coarse additive remainder.

Finally weaken the already obtained support, size and length prices to
the two sets required for the finite-reader comparison. Such weakening
only decreases the positive weights in the rooted sums. The construction
for an arbitrary finite marked family is consequently available for
both sets of prices on this same step, with the same density identity
and the same original partition denominator. Combining these steps with
the common reader and finite-menu choices in
Section~\ref{sec:actual-rank-purchases} supplies its local compiler input.
The callback uses the original history and its binding, as established
in Section~\ref{sec:public-run-binding}.

This application discharges the local matrix, two-inventory,
first-derivative determinant, owner and all-palette premises of the
corrected compiler at the specified fixed patch data. It retains the
selected history's actual path/clip, alongside, inherited-law and
remainder hypotheses until their history-level suppliers are applied.
In particular, it does not assert existence of a history from local
coupling inequalities alone.
